% Part: sets-functions-relations
% Chapter: size-of-sets-complete

\documentclass[../../../include/open-logic-chapter]{subfiles}

\begin{document}

\olchapter{sfr}{siz}{集合的大小}

\begin{editorial}
本章讨论枚举、可数性与不可数性。
其中几节提供了两个版本：一个较初等的版本，将枚举视为列表，或从 $\PosInt$ 出发的满射函数；另一个较抽象的版本，则将枚举定义为与 $\Nat$ 的双射。
\end{editorial}

\olimport{introduction}

\olimport{enumerability}

\olimport{zig-zag}

\olimport{pairing}

\olimport{pairing-alt}

\olimport{non-enumerability}

\olimport{reduction}

\olimport{equinumerous-sets}

\olimport{comparing-size}

\olimport{schroder-bernstein}

\begin{editorial}
以下 \olref[sfr][siz][enm-alt]{sec}、
\olref[sfr][siz][nen-alt]{sec}、\olref[sfr][siz][red-alt]{sec} 是
\olref[sfr][siz][enm]{sec}、\olref[sfr][siz][nen]{sec}、\olref[sfr][siz][red]{sec}
的替代版本，由 Tim Button 为其《Open Set Theory》教材编写。它们
略为进阶，并使用了一种在集合论语境下更合适的不同的可数性定义
（即与 $\Nat$ 或其某个初始段双射，而非是可列表的，或是从 $\PosInt$ 出发的满射函数的值域）。
\end{editorial}

\olimport{enumerability-alt}
\olimport{non-enumerability-alt}
\olimport{reduction-alt}

\OLEndChapterHook

\end{document}